\chapter{Quantas entradas são precisas}
% versão completa: chapters/19_dendrites.tex

O neurônio tem prolongamentos de entrada ramificados: os dendritos. Em alguns neurônios, são uma árvore inteira com milhares de contatos; em outros, um ou dois ramos, ou nada. Para o engenheiro, o número de entradas é um parâmetro do circuito, e ele é escolhido de acordo com a tarefa.

\section*{De zero a cem mil}

Os sensores de toque e de dor não têm dendrito nenhum: o fio simplesmente vai da pele até a medula espinhal, e o sensor fica na ponta dele. É uma linha com o sensor na extremidade distante.

\pencilfig{19}{\pencilart{19}}{Muitas entradas: um somador; uma única entrada enorme: um repetidor rápido e preciso.}

Os neurônios do olfato e os que transmitem o sinal dos sensores do olho e do ouvido têm um dendrito. Uma entrada, uma saída: um repetidor, um buffer.

Os detectores da direção do som têm dois dendritos, voltados para lados opostos, um para cada ouvido. Separando as duas entradas em ramos diferentes, o neurônio torna o ``e'' mais rigoroso: duas porções de corrente no mesmo ramo atrapalhariam uma à outra.

Os neurônios pequenos do circuito de aprendizado motor têm uns quatro dendritos curtos, cada um agarrando a sua entrada. Esses neurônios são dezenas de bilhões, e cada um soma apenas quatro sinais. Em compensação, juntos eles decompõem a entrada em milhões de combinações, e um grande neurônio de saída com cem mil entradas escolhe entre elas as que servem.

E, por fim, os neurônios com árvores enormes e milhares de entradas são somadores e classificadores, em que ramos isolados funcionam quase como pequenos circuitos independentes.

\section*{Por que é assim}

Tudo se decide pela carga. Um neurônio pequeno, com poucas entradas, se carrega com uma única conexão grande em frações de milissegundo: esse neurônio é preciso no tempo e transmite o sinal adiante com segurança. Uma árvore grande nunca vai se carregar com uma conexão só --- a carga vaza antes. Ela precisa de dezenas de sinais chegando juntos, mas em troca pode somá-los e compará-los. Poucas entradas e conexões grandes: para o tempo exato. Muitas entradas: para a contagem.

\pencilfig{128}{%
% charging: a small neuron reaches threshold from one large synapse at once; a big tree barely moves
% from one input and needs many inputs arriving together
\foreach \dx/\t in {0/{poucas entradas}, 4.3/{árvore grande}}{
  \draw[pen] (\dx,0) -- (\dx+3.6,0);
  \draw[pcGray, densely dashed, line width=0.5pt] (\dx,0.9) -- (\dx+3.6,0.9);
  \node[font=\small\itshape, text=pcGray, anchor=south] at (\dx+1.8,1.75) {\t};}
\node[lbl, anchor=east] at (-0.05,0.9) {limiar};
% small neuron
\draw[pcBlue, line width=2pt] (0.6,-0.15) -- (0.6,-0.6);
\draw[penlite=pcBlue] (0,0) -- (0.6,0) -- plot[domain=0.6:0.8, samples=12] (\x,{1.3*(1-exp(-(\x-0.6)/0.18))}) -- (0.8,1.6) -- (0.84,0) -- (3.5,0);
\node[lbl, text=pcRed, anchor=west] at (0.85,1.45) {clique imediato};
\node[lbl, text=pcBlue, anchor=north] at (0.6,-0.62) {uma entrada grande};
% big tree
\draw[pcBlue, line width=0.6pt] (4.8,-0.15) -- (4.8,-0.45);
\foreach \s in {6.15,6.2,6.25,6.3,6.35,6.4,6.45,6.5}{\draw[pcBlue, line width=0.6pt] (\s,-0.15) -- (\s,-0.45);}
\draw[penlite=pcBlue] (4.3,0) -- (4.8,0) -- plot[domain=4.8:6.15, samples=20] (\x,{0.16*((\x-4.8)/0.15)*exp(1-(\x-4.8)/0.15)}) -- plot[domain=6.15:6.62, samples=14] (\x,{1.25*(1-exp(-(\x-6.15)/0.4))}) -- (6.62,1.6) -- (6.66,0) -- (7.9,0);
\node[lbl, anchor=south] at (4.95,0.2) {quase nada};
\node[lbl, text=pcBlue, anchor=north] at (6.33,-0.47) {muitas juntas};
}{Um neurônio pequeno se carrega com uma única conexão grande em frações de milissegundo e clica na hora. Uma árvore grande quase não muda com uma entrada só: precisa de dezenas de sinais chegando juntos.}

Assim chegamos à terceira separação dos neurônios: pela substância, pelo aparelho e pelo número de entradas. Cada uma vê a mesma máquina por um novo ângulo.

\fullversion{19}
